\section{数据来源与预处理}

\subsection{TCGA 与 UCSC Xena 数据资源}

TCGA（The Cancer Genome Atlas）是癌症研究领域最重要的公共多组学数据库之一，覆盖了大量癌种的基因组、转录组、表观组、蛋白组以及临床信息。其数据规模大、注释较完整、样本来源统一，因此非常适合用于多组学方法研究。UCSC Xena 则提供了 TCGA 及其他公共项目的统一下载与可视化入口，使研究者能够较为方便地获取不同组学的标准化矩阵和临床标签。

本研究采用的 TCGA-BRCA 数据包含三类核心信息：RNA-seq 表达矩阵、DNA 甲基化矩阵以及临床亚型标签。当前仓库中已经整理好的数据文件分别位于 \texttt{datasets/TCGA-BRCA.star\_counts.csv}、\texttt{datasets/TCGA-BRCA.methylation450.csv} 与 \texttt{datasets/brca\_pam50.csv}。其中，RNA 和甲基化矩阵均以“特征 × 样本”的形式组织，临床文件用于提供监督学习标签。

\subsection{RNA-seq 数据特点}

RNA-seq 数据在本研究中承担最基础的功能表达信息角色。转录组矩阵记录了每个基因在不同样本中的表达水平，其维度通常较高，且表达分布具有明显偏态：少量高表达基因占据较大信号，而大量低表达基因对分类的贡献有限。因此，RNA 预处理通常需要进行对数变换、低表达过滤与高方差特征筛选。

当前实现采用 \texttt{log2(TPM + 1)} 形式进行数值变换，以压缩极端值并缓解表达量分布不均问题。随后使用均值过滤与方差排序保留高区分度特征。这个步骤在实验中起到两个作用：一是降低噪声，二是避免模型被大量无关基因拖累。对于高维小样本场景而言，这类特征筛选往往比复杂的深度模型更稳定。

\subsection{DNA 甲基化数据特点}

甲基化数据反映的是 CpG 位点上的表观遗传状态，其维度更高，典型平台（如 Illumina 450K）可能包含数十万条 probe。与 RNA 不同，甲基化数据的解释通常更偏向调控层面：某些启动子区域的甲基化升高可能导致相关基因沉默，而不同亚型在甲基化谱上的差异往往反映其调控机制的变化。

由于甲基化数据极高维且稀疏，若直接进入分类器，很容易出现严重过拟合。因此本研究同样采用高方差筛选策略，只保留最具区分度的 probe。当前代码中，甲基化预处理会先删除缺失值，再按方差排序保留 Top-K 特征。该策略虽然简单，但在当前中小样本实验中具有较高的可操作性。

\subsection{临床标签与样本对齐}

多组学分类任务最容易出错的环节不是模型，而是样本对齐。RNA、甲基化与临床标签文件必须对应同一批样本，否则模型学到的只是“错位后的噪声关系”。当前仓库中临床文件采用 TSV 格式，且存在重复样本和标签命名不完全一致的问题，因此在数据加载阶段需要做格式统一：一方面将 \texttt{PAM50} 映射为 \texttt{label}，另一方面将 \texttt{Her2} 统一为 \texttt{HER2}。

样本对齐时，当前实现采用三者样本交集，并保持 RNA 列顺序作为最终对齐顺序，从而保证多个组学在同一个样本索引上进行建模。这个设计避免了集合交集带来的无序问题，也减少了后续拼接或 MOFA 输入阶段的错位风险。

\subsection{训练-测试划分与防泄露原则}

在机器学习任务中，最常见但也最隐蔽的错误之一是信息泄露。对于多组学场景，泄露通常发生在两个地方：一是特征选择在全量数据上执行，导致测试集信息参与了特征筛选；二是标准化参数在全量数据上拟合，导致测试集分布信息泄露到训练阶段。

为降低这种风险，当前实验流程采用先切分再拟合预处理器的策略。具体来说，先对样本进行训练/测试切分，再在训练集上学习标准化参数，并将同一预处理器应用到测试集。对于特征筛选，由于当前实现需要保证训练与测试使用同一组特征，因此本仓库采用先在全量数据上确定特征集合、再在训练集上拟合标准化器的方案。这一做法兼顾了特征维度一致性与训练阶段的统计拟合稳定性，但在严格意义上仍需要在论文中明确其方法边界。

\subsection{标签映射}

当前实验仅保留四个主要 PAM50 亚型，并采用如下数值映射：
\begin{itemize}
  \item LumA $\rightarrow$ 0
  \item LumB $\rightarrow$ 1
  \item HER2 $\rightarrow$ 2
  \item Basal $\rightarrow$ 3
\end{itemize}

这种编码方式使监督学习模型能够直接处理分类目标。同时，在报告结果时仍建议保留原始亚型名称与数值标签的对应说明，以便读者理解分类结果。
